Para a minimização da função $w_{13}$ definida para o trabalho, o algoritmo da genético com restrição no espaço de busca tem a implementação um pouco mais difícil pois precisa manipular cadeias de bits, porém dispõe da mesma capacidade de contornar o problema da não-convexidade da função $w_{13}$, que possui vários mínimos locais, devido a seu caráter estocástico. O algoritmo genético funciona através do cruzamento e mutação dos elementos das cadeias de bits da população de forma semi-aleatória. Determina-se, através de uma função de aptidão que deve ser sempre positiva, o valor de aptidão atingido por uma determinada cadeia de bits. Após o registro destas variáveis, é feito o cálculo da chance de cruzamento que cada cadeia dispõe em relação às outras, membros com aptidão maior tem mais chance de serem selecionados. Finalmente, uma pequena taxa de mutação é adicionada a cada bit de forma a garantir variedade populacional. 